\section{用量词精确刻画数列趋近}
\label{sec:12-Real-Analysis-Support/02-Convergence/01}

微积分入门课里说``\(n\) 很大时 \(x_n\) 接近 \(L\)''。这句话抓到了直觉，但用来做证明不够。细看两个反例。

反例一：\(x_n=(-1)^n/n\)。它反复穿过 0，每次穿过时离 0 很近，但随即又远离。它最终趋近 0 吗？直觉说``是的，因为振幅在衰减''——但你需要说清``反复路过附近又离开''为什么仍然算趋近。

反例二：\(x_n=1+1/n\)。它单调递减，越来越接近 0——每一步都比上一步更靠近 0。但它的极限是 1，不是 0。``越来越接近''跟``最终停在附近''是两回事。只在 1 附近最终停留。

两个反例指向同一个缺口：日常语言里的``接近''分不出``暂时路过''和``永久停留''，也分不出``朝某个方向移动''和``朝某个位置收敛''。你需要一个能一锤定音的定义。

\subsection{\(\varepsilon\)--\(N\) 定义：把``最终停留''变成可操作的量词}

数列 \((x_n)\) 收敛到 \(L\)，记作 \(x_n\to L\)，如果

\[
\forall\varepsilon>0,\ \exists N\in\N,\ \forall n\ge N,\ |x_n-L|<\varepsilon.
\]

从左往右逐块翻译：

\begin{itemize}
\tightlist
\item \(\forall\varepsilon>0\)：无论你要求多小的误差容限；
\item \(\exists N\in\N\)：我都能找到一个尾部起点（序号 \(N\)）；
\item \(\forall n\ge N\)：从这个起点之后的每一项；
\item \(|x_n-L|<\varepsilon\)：都落在以 \(L\) 为中心、\(\varepsilon\) 为半径的带子里。
\end{itemize}

量词顺序决定了一切。\(\varepsilon\) 在前，所以 \(N\) 可以依赖 \(\varepsilon\)——误差要求越严，等待时间可能越长。但 \(N\) 一旦选定，就要同时压制它之后的所有项。不存在``第 100 项满足第 101 项不满足，第 102 项又满足''——一旦进入带子，就不能再出来。

回头检验两个反例。对 \((-1)^n/n\)：给定 \(\varepsilon=0.1\)，取 \(N=11\)，因为 \(|(-1)^n/n|=1/n\le 1/11<0.1\) 对所有 \(n\ge 11\) 成立。无论正负，振幅被 \(1/n\) 压住了。对 \(1+1/n\) 和 0 的关系：给定 \(\varepsilon=0.5\)，不存在 \(N\) 使 \(n\ge N\) 时 \(|1+1/n-0|<0.5\) 总是成立——因为对任何 \(n\ge1\)，\(1+1/n\) 都大于 1，距离至少是 1，根本进不了 0.5 的带子。所以它不收敛到 0。

\subsection{一个标准证明：先算误差再反解 \(N\)}

证明

\[
\frac{3n+1}{n+2}\to 3.
\]

直接计算误差：

\[
\left|\frac{3n+1}{n+2}-3\right|
= \left|\frac{(3n+1)-3(n+2)}{n+2}\right|
= \left|\frac{-5}{n+2}\right|
= \frac{5}{n+2}.
\]

目标：使这个量小于给定的 \(\varepsilon\)。解不等式

\[
\frac{5}{n+2} < \varepsilon
\iff n > \frac{5}{\varepsilon}-2.
\]

只要取

\[
N > \max\!\left\{0,\ \frac{5}{\varepsilon}-2\right\},
\]

则对所有 \(n\ge N\)，误差小于 \(\varepsilon\)。\(\max\) 里塞一个 0 只是防止 \(\varepsilon\) 很大的时候算出负数——这种防御性书写在证明中常见但无关要害。

证明的要领是：\textbf{先把目标误差改写成一个可以用 \(n\) 控制的上界，再从这个上界反解出 \(N\)}。不要一上来就猜一个 \(N\) 然后验证——那只有在极其简单的例子里碰巧行得通。

\subsection{极限唯一：用三角不等式逼出矛盾}

若 \(x_n\to L\) 同时 \(x_n\to M\)，且 \(L\ne M\)。取

\[
\varepsilon = \frac{|L-M|}{3}.
\]

由于 \(x_n\to L\)，存在 \(N_1\) 使 \(n\ge N_1\) 时 \(|x_n-L|<\varepsilon\)。由于 \(x_n\to M\)，存在 \(N_2\) 使 \(n\ge N_2\) 时 \(|x_n-M|<\varepsilon\)。取 \(N=\max\{N_1,N_2\}\)，则对 \(n\ge N\)，两个不等式同时成立。三角不等式给出

\[
|L-M| \le |L-x_n|+|x_n-M| < \varepsilon+\varepsilon = 2\varepsilon = \frac{2}{3}|L-M|,
\]

即 \(|L-M|<\frac{2}{3}|L-M|\)，矛盾。所以 \(L=M\)。

选 \(\varepsilon = |L-M|/3\) 的用意：如果两个极限不同，它们之间有某个正距离。把这个距离的三分之一当作容许误差，那么同一个 \(x_n\) 不可能同时离 \(L\) 和 \(M\) 都这么近——它必须选一边。参数 \(1/3\) 的精确数值不重要，小于 \(1/2\) 就行，只是 \(1/3\) 让代数干净。

\subsection{收敛数列有界：尾部被极限框住，前面只有有限个}

若 \(x_n\to L\)，取 \(\varepsilon=1\)，由收敛定义存在 \(N\) 使 \(n\ge N\) 时

\[
|x_n-L|<1 \implies L-1 < x_n < L+1.
\]

尾部有界了。前面 \(x_1,\ldots,x_{N-1}\) 只有有限项，取它们绝对值的最大值与 \(|L|+1\) 比较，全序列有界。

有界不推收敛——\((-1)^n\) 就是个钉子户。但它逼近了一个更精细的事实：有界序列虽然可能不收敛，但一定有收敛的子列。

\subsection{Bolzano--Weierstrass：有界必有收敛子列}

这个定理说：每个有界实数序列都含有收敛子列。证明的骨架不涉及极限的构造，而是巧妙地利用单调性。

先从任意序列中提取单调子列。把一项 \(x_k\) 称作风上的``峰''，如果它之后的所有项都不超过它：对所有 \(m>k\)，\(x_m\le x_k\)。

\begin{itemize}
\tightlist
\item 若有无穷多个峰：这些峰按序号递增排列，但由于每个峰压住它后面的所有项，排在后面的峰小于等于排在前面的峰——它们构成递减子列。
\item 若只有有限个峰：取最后一个峰之后的下一项作为起点。它不是峰，所以后面存在比它大的项；取那个更大的项，它也不是峰（因为我们已经过了最后一个峰），后面又存在比它更大的项；继续下去，得到一个严格递增子列。
\end{itemize}

无论哪种情形，都得到一个单调子列。若原序列有界，这个单调子列自然也有界。上一节证明了单调有界必收敛，于是该子列收敛。证毕。

Bolzano--Weierstrass 依赖于实数完备性——单调有界收敛定理是完备性的推论，而提取单调子列的技巧不依赖完备性，在任何全序集上都成立。在 \(\R^d\) 中定理同样成立（逐分量提取子列）。下一单元会把它识别为紧致性的序列表述——有界闭集上的每个序列都有收敛子列。

\subsection{子列能暴露不收敛}

若 \(x_n\to L\)，它的任何子列也收敛到同一个 \(L\)。因此，只要能找到两个子列趋向不同极限（或一个发散的子列），就能断定原序列不收敛。

对 \(x_n=(-1)^n\)：偶数子列恒为 1，奇数子列恒为 \(-1\)。两个极限不同，原序列不收敛。

反过来不成立：有一个收敛子列不能说明原序列收敛。每个有界序列都有收敛子列（Bolzano--Weierstrass），但完全可以在多个聚点之间来回漂移。

\subsection{单调有界：最简单的收敛保证}

若 \((x_n)\) 单调递增且有上界，则

\[
x_n \to \sup\{\,x_n:n\in\N\,\}.
\]

单调递减且有下界时极限为下确界。这三件事各自不可或缺：

\begin{itemize}
\tightlist
\item 单调性防止振荡——序列只能朝一个方向走；
\item 有界性防止逃逸——它不能走到无限远；
\item 完备性提供着陆点——那个上确界必须存在且属于这个空间。
\end{itemize}

三条拆开检验：\(x_n=n\) 单调但无上界，发散；\((-1)^n\) 有界但不单调，振荡；在 \(\Q\) 中取递增有理数列逼近 \(\sqrt2\)，单调且有有理上界，但空间不完备，无极限。

很多迭代算法的收敛证明都在复现这个模板。先证序列单调（通常借助目标函数的单调递减性质），再构造一个下界（能量函数或损失函数的下界），于是极限存在。但至此只证明了``往某个地方走且不会走过头''——极限满足什么方程需要额外的工作。

\subsection{代数运算与夹逼：收敛性质的传播}

若 \(x_n\to x\)，\(y_n\to y\)，则：

\[
x_n+y_n\to x+y,\qquad x_ny_n\to xy,
\]

且若 \(y\ne0\) 且尾部项不为零，则 \(x_n/y_n \to x/y\)。

这些规则不应沦为``把极限符号当作普通代数操作''的口诀，而要看到背后的误差估计机制。以加法为例：给定 \(\varepsilon>0\)，分别对两个数列取 \(N_1,N_2\) 使各自误差小于 \(\varepsilon/2\)，取 \(N=\max\{N_1,N_2\}\)，则对 \(n\ge N\)，

\[
|(x_n+y_n)-(x+y)| \le |x_n-x|+|y_n-y| < \frac{\varepsilon}{2}+\frac{\varepsilon}{2}=\varepsilon.
\]

乘积多一步：\(|x_ny_n-xy| \le |x_n||y_n-y| + |y||x_n-x|\)，其中 \(|x_n|\) 有界（收敛数列有界），所以第二项由 \(|x_n-x|\) 控制，第一项由 \(|y_n-y|\) 控制。商的证明更精细：先用 \(|y_n|>|y|/2\)（从 \(|y_n-y|<|y|/2\) 推出，对充分大的 \(n\) 成立）把分母控制在一个正常数以下。

夹逼定理的证明只有一行量词：若

\[
a_n \le x_n \le b_n
\]

且 \(a_n,b_n\to L\)，则 \(x_n\to L\)。证明——给定 \(\varepsilon>0\)，取 \(N\) 使 \(n\ge N\) 时同时有 \(L-\varepsilon<a_n\) 和 \(b_n<L+\varepsilon\)，于是

\[
L-\varepsilon < a_n \le x_n \le b_n < L+\varepsilon,
\]

即 \(|x_n-L|<\varepsilon\)。关键条件是上下两边夹向同一个极限；只知两边有界远远不够。
